\chapter{Fine-Tuning Agents for Domain Behaviour}

\begin{quote}
\textit{``A legal tech company deployed GPT-4o to extract structured data from contracts. The model was excellent at understanding English, but consistently misread their internal contract taxonomy — confusing 'Effective Date' with 'Execution Date,' and 'Counterparty' with 'Licensor.' After three months of prompt engineering with minimal improvement, they switched strategy: they collected 800 correct extraction examples from their own paralegals, trained a fine-tuned version of Llama-3-8B using DPO, and deployed it. The fine-tuned 8B model outperformed GPT-4o on their internal taxonomy by 23 percentage points, and cost 96\% less per inference.''}
\end{quote}

Fine-tuning is not about making a model ``smarter'' in general. It is about encoding domain-specific knowledge and behavioral preferences that cannot be expressed efficiently in a prompt. This chapter covers the three techniques that matter in production: DPO for behavioral alignment, LoRA for efficient training, and synthetic data generation for building training sets.

\section{When to Fine-Tune vs. Prompt Engineer}

This is the most important decision in this chapter. Fine-tuning has a significant upfront cost (data collection, training, infrastructure). Prompt engineering has near-zero cost but limited ceiling.

\begin{center}
\begin{tabular}{|l|c|c|}
\hline
\textbf{Situation} & \textbf{Prompt Engineering} & \textbf{Fine-Tuning} \\
\hline
General capability (reasoning, writing) & \ding{51} & \ding{55} \\
Domain-specific vocabulary/taxonomy & Partial & \ding{51} \\
Consistent output format/schema & Partial & \ding{51} \\
Reducing inference cost (smaller model) & \ding{55} & \ding{51} \\
Specific behavioral style (tone, persona) & Partial & \ding{51} \\
Frequent updates (daily schema changes) & \ding{51} & \ding{55} \\
\hline
\end{tabular}
\end{center}

\textbf{Rule of thumb:} If you have $>$500 examples of preferred behavior and the behavior is stable, fine-tuning will outperform any prompt.

\section{Direct Preference Optimization (DPO)}

Standard supervised fine-tuning (SFT) trains the model to replicate correct outputs. DPO is different: it trains the model using \textit{preference pairs} — examples of correct and incorrect behavior — and directly optimizes the model to \textit{prefer} the correct behavior.

\subsection{The DPO Loss Function}
Given a prompt $x$, a preferred response $y_w$ (``winner''), and a rejected response $y_l$ (``loser''):
\begin{equation}
\mathcal{L}_{\text{DPO}}(\pi_\theta; \pi_{\text{ref}}) = -\mathbb{E}_{(x, y_w, y_l) \sim \mathcal{D}} \left[ \log \sigma \!\left( \beta \log \frac{\pi_\theta(y_w \mid x)}{\pi_{\text{ref}}(y_w \mid x)} - \beta \log \frac{\pi_\theta(y_l \mid x)}{\pi_{\text{ref}}(y_l \mid x)} \right) \right]
\end{equation}

where:
\begin{itemize}
    \item $\pi_\theta$ is the policy (model being trained)
    \item $\pi_{\text{ref}}$ is the frozen reference model (the pre-trained base)
    \item $\beta$ is the temperature parameter controlling how strongly to diverge from the reference (typically 0.1 -- 0.5)
    \item $\sigma$ is the sigmoid function
\end{itemize}

\textbf{Intuition:} The loss increases the log-probability of $y_w$ relative to how the reference model rated it, while simultaneously decreasing the log-probability of $y_l$ — all while a KL penalty (implicit in the ratio terms) prevents the model from diverging catastrophically from the base.

\subsection{Building a Tool-Call DPO Dataset}
For agent fine-tuning, preference pairs are tool-call trajectories:


\begin{center}
\noindent\rule{0.6\textwidth}{0.4pt} \\
\vspace{0.3em}
\textit{[Code implementation is available in the companion coding handbook: \texttt{coding-handbook/ch17\_fine\_tuning/}]} \\
\noindent\rule{0.6\textwidth}{0.4pt}
\end{center}


\section{LoRA: Training Only What Matters}

A 70B parameter model cannot be fine-tuned on a single GPU — the weights alone require 140GB of VRAM in FP16. LoRA (Low-Rank Adaptation) sidesteps this by \textit{freezing} all original weights and training only small low-rank adapter matrices that are added to specific layers.

\subsection{LoRA Mathematics}
For a pre-trained weight matrix $W_0 \in \mathbb{R}^{d \times k}$, LoRA adds a low-rank perturbation:
\begin{equation}
W = W_0 + \Delta W = W_0 + BA
\end{equation}
where $B \in \mathbb{R}^{d \times r}$, $A \in \mathbb{R}^{r \times k}$, and rank $r \ll \min(d, k)$.

\textbf{Parameter count comparison} for a single attention layer of Llama-3-8B ($d=4096$, $k=4096$):
\begin{align}
\text{Full fine-tune:} &\quad 4096 \times 4096 = 16{,}777{,}216 \text{ parameters} \\
\text{LoRA (r=16):} &\quad 4096 \times 16 + 16 \times 4096 = 131{,}072 \text{ parameters}
\end{align}

LoRA trains \textbf{128$\times$ fewer parameters} while achieving comparable performance for domain adaptation tasks.


\begin{center}
\noindent\rule{0.6\textwidth}{0.4pt} \\
\vspace{0.3em}
\textit{[Code implementation is available in the companion coding handbook: \texttt{coding-handbook/ch17\_fine\_tuning/}]} \\
\noindent\rule{0.6\textwidth}{0.4pt}
\end{center}


\section{Synthetic Data Generation with a Teacher Model}

Collecting 500+ high-quality preference pairs from human experts is expensive. A faster alternative is using a strong ``teacher'' model (GPT-4o) to generate training data for a weaker ``student'' model (Llama-3-8B).


\begin{center}
\noindent\rule{0.6\textwidth}{0.4pt} \\
\vspace{0.3em}
\textit{[Code implementation is available in the companion coding handbook: \texttt{coding-handbook/ch17\_fine\_tuning/}]} \\
\noindent\rule{0.6\textwidth}{0.4pt}
\end{center}

